\section{An exponential bound for order-reversing images}
\label{amp:section}

We now prove Theorem~\ref{amp:image}, combining transport and the
nested-chain inequality to bound order-reversing relation images.
Write $\mathcal R_H$ for the monoid of all binary relations
on a finite set $H$, with multiplication in path order.  For
$F\subseteq H$, write $I_F=\{(x,x):x\in F\}$; relations on $F$ are
also regarded as relations on $H$ supported on $F\times F$.

The proof amplifies one off-diagonal pair in the image into a family
of pair additions.  The chain inequality bounds their total cost in
missing recurrent classes; restriction to smaller relation monoids
bounds each cost from below.  We first record two elementary facts
that allow these restrictions without losing the hypotheses.
The pattern of relation additions follows the matching-diagram
argument in \cite[Section~3]{CompanionDeterminization2026}; all loss estimates
needed here are supplied by the path-diagram results just proved.

\subsection{Corners and unit lifts}

For an idempotent $d$ in a semigroup $A$, its \emph{corner}
$dAd=\{dad:a\in A\}$ is a monoid with identity $d$.
A \emph{unit} means an element with a two-sided inverse relative to
the specified monoid identity.

\begin{lemma}[Minimal idempotent corners]
\label{amp:minimal-corner}
Let $A$ be a finite semigroup, let $M$ be a monoid, and let
$\theta:A\to M$ be a surjective multiplicative map.  For every
idempotent $q\in M$, there is an idempotent $d\in A$ with
$\theta(d)=q$ that is minimal among such idempotents under
$\trianglelefteq$.  For every such minimal $d$, restriction of
$\theta$ gives a unital surjection
\[
 dAd\longrightarrow qMq.
\]
Every element of $dAd$ whose image is a unit of $qMq$ is itself a
unit of $dAd$.  In particular, every unit of $qMq$ has a unit lift.
\end{lemma}

\begin{proof}
Every element $w$ of a finite semigroup has an idempotent positive
power: the sequence of powers is eventually periodic, so a
sufficiently large exponent divisible by its eventual period gives
$w^{2k}=w^k$.  Applying this to a preimage of $q$ gives an
idempotent preimage of $q$.  Finiteness then gives one minimal in
the nesting order.

For any $y\in qMq$, choose $x\in A$ with $\theta(x)=y$.
Then $\theta(dxd)=qyq=y$, proving surjectivity on the corner;
its identity $d$ maps to $q$.  Now let $w\in dAd$ have unit image.
An idempotent power $w^k$ has an image that is both a unit and an
idempotent in $qMq$, hence equals $q$.  Since
$w^k\trianglelefteq d$, minimality gives $w^k=d$.  If $k>1$,
the element $w^{k-1}$ is a two-sided inverse of $w$ in $dAd$;
if $k=1$, then $w=d$ is already the identity.
\end{proof}

The next lemma identifies a smaller full relation monoid inside an
idempotent corner.  Its hypothesis says that the old relation
restricts to the identity on the retained points.

\begin{lemma}[Restriction of a relation corner]
\label{amp:relation-corner}
Let $P\in\mathcal R_H$ be idempotent, let $F\subseteq H$, and
put $i=I_F$.  Suppose that $iPi=i$, and define $Q=PiP$.
Then $Q$ is idempotent, $Q\trianglelefteq P$, and restriction is
a unital monoid isomorphism
\begin{equation}
 \rho:Q\mathcal R_HQ\longrightarrow\mathcal R_F,
 \qquad \rho(Z)=iZi.
 \label{amp:restriction}
\end{equation}
Its inverse sends $A\in\mathcal R_F$ to $PAP$.
\end{lemma}

\begin{proof}
Using $P^2=P$ and $iPi=i$, we obtain
\[
 Q^2=Q,\qquad PQ=QP=Q,\qquad
 Qi=Pi,\qquad iQ=iP,\qquad iQi=i.
\]
If $Z=QZQ$, then $PZP=Z$, and therefore
\[
 Z=PiPZPiP=PiZiP.
\]
Thus restriction determines $Z$ uniquely.  Conversely, for a
relation $A=iAi$ supported on $F$, the relation $PAP$ belongs to
$Q\mathcal R_HQ$, since $PAP=QAQ$, and
\[
 i(PAP)i=(iPi)A(iPi)=A.
\]
This proves bijectivity with the claimed inverse.  For supported
relations $A,B\in\mathcal R_F$,
\[
 (PAP)(PBP)=PAPBP
           =PA(iPi)BP
           =PABP.
\]
Hence the inverse, and therefore restriction, preserves
multiplication.  Finally $\rho(Q)=i$, the identity of
$\mathcal R_F$.
\end{proof}

\subsection{Amplifying a single added pair}

Set
\begin{equation}
 t=64,\qquad
 D(h)=2^{\lfloor(h-2)/(2t-1)\rfloor}\quad(h\geq2).
 \label{amp:D}
\end{equation}
The choice $t=64$ balances the two estimates proved below: $t^2$
pair additions each cost at least half the smaller-instance bound,
while their total cost is at most $16t$ times the original loss.
The resulting gain is $t/32=2$, at a reduction of $2t-1=127$ points.
We prove a missing-set bound strong enough to survive passage to
any of the corners above.  The additional unit-lifting hypothesis
permits conjugation by arbitrary permutations of $H$; Lemma
\ref{amp:minimal-corner} will supply it when we return to Theorem
\ref{amp:image}.

\begin{proposition}[Cost of one added pair]
\label{amp:single-edge}
Let $m\geq1$, let $e\in\mathcal T_m$ be idempotent, and let
$S\subseteq e\mathcal T_me$ be a submonoid with identity $e$.
Let $H$ be a finite set of cardinality $h\geq2$.  Suppose that
$\phi:S\to\mathcal R_H$ is a unital surjective homomorphism
such that
\[
 z\sqsubseteq w\ \Longrightarrow\ \phi(w)\subseteq\phi(z)
 \qquad(z,w\in S),
\]
and suppose that each permutation of $H$ has a lift that is a
unit of $S$.  If $a\in S$ is idempotent and, for distinct $x,y\in H$,
\[
 \phi(a)=I_H\cup\{(x,y)\},
\]
then $|\mathcal M_e(a)|\geq D(h)$.
\end{proposition}

\begin{proof}
We induct on $h$, simultaneously over all the data in the
statement, including $m$, $e$, $S$, and $\phi$.
For $2\leq h\leq2t$, we have $D(h)=1$.
If $\mathcal M_e(a)$ were empty, Lemma
\ref{diag:rectangles} would give $e\sqsubseteq a$.
Order reversal would then give
$I_H\cup\{(x,y)\}\subseteq I_H$, a contradiction.

Assume henceforth that $h\geq2t+1$, and put
$c=|\mathcal M_e(a)|$.  We will construct $t^2$ nested steps
with total missing-set size at most $16tc$.  Each step will
contain a smaller instance of the Proposition, on $h-2t+1$
points.

\medskip
\noindent\textbf{Step 1: A common budget for all pair additions.}
Choose disjoint sets
\[
 U=\{u_1,\ldots,u_t\},\qquad
 V=\{v_1,\ldots,v_t\},\qquad \{z\}
\]
in $H$.  Conjugating $a$ by unit lifts of permutations gives
elements $s_i,s'_j\in S$ with
\begin{equation}
 \phi(s_i)=I_H\cup\{(u_i,z)\},\qquad
 \phi(s'_j)=I_H\cup\{(z,v_j)\}.
 \label{amp:hub}
\end{equation}
Indeed, permutation conjugation acts transitively on ordered
pairs of distinct points, and the inverse of a unit lift maps
to the inverse permutation.  Explicitly, a unit $p\in S$
gives the conjugate $pap^{-1}$, with $pep^{-1}=e$.
Applying Lemma \ref{diag:transport} with
$J=\mathcal M_e(a)$ therefore gives
\[
 |\mathcal M_e(s_i)|\leq2c,\qquad
 |\mathcal M_e(s'_j)|\leq2c.
\]
Consequently the single set
\[
 J=\bigcup_{i=1}^t\mathcal M_e(s_i)
       \ \cup\ \bigcup_{j=1}^t\mathcal M_e(s'_j)
       \ \subseteq\mathcal C_e
\]
has $|J|\leq4tc$, and the product inequality of Lemma
\ref{diag:missing} gives
$\mathcal M_e(s_i s'_j)\subseteq J$ for every $i,j$.

Let $P_0=I_{H\setminus\{z\}}$.  Choose an idempotent lift
$b_0\in S$ of $P_0$, by taking an idempotent positive power
of any lift, and define
\[
 g_{ij}=b_0s_i s'_j b_0\in b_0Sb_0.
\]
The image of $s_i s'_j$ consists of the identity and the three
pairs $(u_i,z)$, $(z,v_j)$, and $(u_i,v_j)$.
The outer factors $P_0$ remove the first two, so
\begin{equation}
 \phi(g_{ij})=P_0\cup\{(u_i,v_j)\}.
 \label{amp:pair-image}
\end{equation}
Figure~\ref{fig:hub} illustrates this conversion of two generator
families into a family indexed by $U\times V$.

\begin{figure}[htbp]
\centering
\begin{tikzpicture}[>=Stealth,font=\small,thick]
 \node (u) at (0,0) {$u_i$};
 \node (z) at (1.8,0) {$z$};
 \node (v) at (3.6,0) {$v_j$};
 \draw[->] (u)--node[above] {$\phi(s_i)$}(z);
 \draw[->] (z)--node[above] {$\phi(s'_j)$}(v);
 \node[align=center] at (1.8,-0.8) {Compose the two relations};
 \draw[->,double] (4.35,0)--node[above] {$P_0\,(\cdot)\,P_0$}(6.35,0);
 \node (uu) at (7.1,0) {$u_i$};
 \node (vv) at (9.5,0) {$v_j$};
 \draw[->] (uu)--node[above] {$\phi(g_{ij})$}(vv);
 \node[align=center] at (8.3,-0.8) {Remove endpoints at $z$};
\end{tikzpicture}
\caption{The relation images of the hub construction, with identity
pairs omitted. Composing the $u_i\to z$ and $z\to v_j$ additions
creates $u_i\to v_j$; sandwiching by $P_0$ removes the two pairs
incident to $z$. The $2t$ generators yield all $t^2$ pair additions.}
\label{fig:hub}
\end{figure}


Since $b_0eb_0=b_0$, the common-set conclusion of Lemma
\ref{diag:transport} gives a \emph{single} set
$J_0\subseteq\mathcal C_{b_0}$ such that
\begin{equation}
 |J_0|\leq8tc,\qquad
 \mathcal M_{b_0}(g_{ij})\subseteq J_0
 \quad(1\leq i,j\leq t).
 \label{amp:common-budget}
\end{equation}

\medskip
\noindent\textbf{Step 2: A nested chain spending that budget.}
List the $t^2$ pairs of $U\times V$ in any order, without
repetitions.  Let $E_\ell$ be the set of the first $\ell$
pairs, and put $P_\ell=P_0\cup E_\ell$ for
$0\leq\ell\leq t^2$.  For any $E,E'\subseteq U\times V$,
\begin{equation}
 (P_0\cup E)(P_0\cup E')=P_0\cup E\cup E'.
 \label{amp:union}
\end{equation}
To check this, $P_0$ acts as the identity on all the endpoints
in question, and no pair of $E$ can be followed by a pair of
$E'$, because $U\cap V=\varnothing$.
In particular, every $P_\ell$ is idempotent.

If the pair at position $\ell$ is $(u_i,v_j)$, choose
$b_\ell$ to be an idempotent positive power of
$b_{\ell-1}g_{ij}b_{\ell-1}$.
Inductively, Equation \eqref{amp:union} gives
$\phi(b_\ell)=P_\ell$.
The sandwich, and hence all its positive powers, is absorbed
on both sides by $b_{\ell-1}$.  Thus
\[
 b_\ell\trianglelefteq b_{\ell-1}
 \qquad(1\leq\ell\leq t^2).
\]
All these elements lie in the $b_0$-corner.  Starting with
$\mathcal M_{b_0}(b_0)=\varnothing$, repeated use of Lemma
\ref{diag:missing} and Equation \eqref{amp:common-budget}
shows that
\[
 \mathcal M_{b_0}(b_\ell)\subseteq J_0
 \qquad(0\leq\ell\leq t^2).
\]
Here taking a positive power cannot enlarge the missing set:
all factors have the same missing set, and their union is
unchanged.  Proposition \ref{nest:chain} now yields
\begin{equation}
 \sum_{\ell=1}^{t^2}\delta_\ell\leq16tc,
 \qquad
 \delta_\ell=|\mathcal M_{b_{\ell-1}}(b_\ell)|.
 \label{amp:total-cost}
\end{equation}
We have bounded the total cost of the chain using only the
$2t$ original conjugates.  It remains to give an inductive
lower bound for each of its $t^2$ steps.

\medskip
\noindent\textbf{Step 3: A smaller instance inside every step.}
Fix $\ell$, let $(u_i,v_j)$ be its newly added pair, and
write $b=b_{\ell-1}$ and $P=P_{\ell-1}$.  Retain the two
endpoints of this pair and all points outside the construction:
\[
 F=\{u_i,v_j\}\cup
       \bigl(H\setminus(U\cup V\cup\{z\})\bigr),
 \qquad i_F=I_F,\qquad Q=Pi_FP.
\]
Then $|F|=h-2t+1\geq2$.  Among the possible extra pairs
of $P$, only $(u_i,v_j)$ has both endpoints in $F$, and this
pair has not yet been added.  Therefore
$i_FPi_F=i_F$.  By Lemma \ref{amp:relation-corner},
$Q\trianglelefteq P$ and restriction gives
$Q\mathcal R_HQ\simeq\mathcal R_F$.
Moreover, Equation \eqref{amp:union} gives
$PP_\ell P=P_\ell$.  Using $i_FQ=i_FP$ and
$Qi_F=Pi_F$, we obtain
\begin{equation}
 i_F(QP_\ell Q)i_F
 =i_FP_\ell i_F
 =I_F\cup\{(u_i,v_j)\}.
 \label{amp:compressed-pair}
\end{equation}

Consider the finite monoid $A=bSb$.  The restriction of
$\phi$ maps it onto $P\mathcal R_HP$: sandwiching any
preimage by $b$ gives a preimage of its $P$-sandwich.
Apply Lemma \ref{amp:minimal-corner} \emph{inside $A$} to
the idempotent $Q$.  We obtain an idempotent $d\in A$ with
$\phi(d)=Q$, minimal there among idempotent lifts of $Q$.
Its corner $dAd$ maps unitally onto
\[
 Q(P\mathcal R_HP)Q=Q\mathcal R_HQ,
\]
and all units of this image have unit lifts.  Composing with
restriction gives a unital surjective homomorphism
\begin{equation}
 \psi:dAd\longrightarrow\mathcal R_F,
 \qquad \psi(w)=i_F\phi(w)i_F.
 \label{amp:smaller-map}
\end{equation}
This map is order-reversing because $\phi$ is and restriction
preserves inclusion.  It has unit lifts of every permutation
of $F$, by the isomorphism in Lemma
\ref{amp:relation-corner}.  Also $dAd$ is a submonoid of
$d\mathcal T_md$ with identity $d$.
These observations verify all the map and domain hypotheses
needed for induction.

Because $b_\ell\trianglelefteq b$, we have $b_\ell\in A$.
By Equation \eqref{amp:compressed-pair}, the element
$db_\ell d\in dAd$ maps under $\psi$ to
$I_F\cup\{(u_i,v_j)\}$.  This relation is idempotent, so
an idempotent positive power $v$ of $db_\ell d$ has the same
image.  Since $d\in A$ is idempotent, $dbd=d$.
Apply Lemma \ref{diag:transport} from the identity $b$ to
the identity $d$, with context $(d,d)$ and missing set
$\mathcal M_b(b_\ell)$.  Then apply the power consequence
of Lemma \ref{diag:missing} in the $d$-corner.  This gives
\begin{equation}
 |\mathcal M_d(v)|
 \leq |\mathcal M_d(db_\ell d)|
 \leq 2\delta_\ell.
 \label{amp:step-cost}
\end{equation}
The induction hypothesis applies to $d$, $dAd$, $\psi$,
and $v$, with the smaller parameter $|F|=h-2t+1$.
It follows that
\[
 2\delta_\ell\geq D(h-2t+1).
\]

Summing this lower bound and using Equation
\eqref{amp:total-cost}, we conclude that
\[
 16tc\ \geq\ \frac{t^2}{2}D(h-2t+1),
 \qquad
 c\ \geq\ \frac{t}{32}D(h-2t+1).
\]
For $t=64$, the last expression is
$2D(h-127)=D(h)$, by the definition of $D$.
This completes the induction.
\end{proof}

\begin{proof}[Proof of Theorem \ref{amp:image}]
Apply Lemma \ref{amp:minimal-corner} to $\phi_0$ and the
idempotent $I_H$.  It supplies an idempotent $e\in S_0$
such that $S=eS_0e$ maps unitally onto $\mathcal R_H$ and
every permutation has a unit lift.  The restricted map
retains order reversal, and $S$ is a submonoid of
$e\mathcal T_me$ with identity $e$.

Choose distinct $x,y\in H$ and any lift in $S$ of
$I_H\cup\{(x,y)\}$.  An idempotent positive power $a$
of that lift has the same image.  Proposition
\ref{amp:single-edge} therefore gives
\[
 D(h)\leq|\mathcal M_e(a)|\leq|\mathcal C_e|\leq2m,
\]
as required.
\end{proof}
